\documentclass[11pt]{article}
\usepackage[margin=1in]{geometry}
\usepackage{amsmath,amssymb,amsthm}
\usepackage[hidelinks]{hyperref}
\newtheorem{theorem}{Theorem}
\newtheorem{lemma}[theorem]{Lemma}
\title{The Narayana polynomial of a ribbon lattice path matroid}
\author{}
\date{25 November 2026}
\begin{document}
\begin{center}
{\LARGE The Narayana polynomial of a ribbon lattice path matroid\par}
\vspace{0.8em}
25 November 2026
\end{center}
\begin{abstract}
The cycle matroid of a polygon with noncrossing diagonals has the same
$Z$-polynomial as its boundary cycle. We derive this observation from the
deletion formula of Braden and Vysogorets. The multi-fan representation of
snake matroids then identifies the $Z$-polynomial of every connected ribbon
lattice path matroid of rank $r$ with the Narayana polynomial of order $r+1$.
This answers a question of Alexandersson.
\end{abstract}

\section{Deleting a diagonal}
Write $P_M$ and $Z_M$ for the Kazhdan--Lusztig and $Z$-polynomials of a
matroid, using the same ranked-lattice convention for a matroid with
loops. In terms of its lattice of flats $L(M)$,
\[
 Z_M(t)=\sum_{F\in L(M)}t^{\mathrm{rk}(F)}P_{M/F}(t).
\]
Both polynomials depend only on the ranked lattice. Write $M(G)$ for the
cycle matroid of a graph. We will use the following immediate consequence
of the deletion formula and the parallel-connection calculation of
Braden and Vysogorets \cite{bv}.

\begin{lemma}\label{lem}
Let a simple graph $G$ be the union of two subgraphs $H_1,H_2$ meeting
precisely in an edge $e=uv$ and its endpoints. If each graph $H_i-e$ is
connected, then $Z_{M(G)}=Z_{M(G-e)}$.
\end{lemma}
\begin{proof}
The cycle matroid $M=M(G)$ is simple, and a $u$-to-$v$ path in either
$H_i-e$ puts $e$ on a cycle. Thus $e$ is a noncoloop and Theorem~2.8,
equation~(6), of \cite{bv} applies. The difference is a sum over flats $F$
such that $e\notin F$ and $F\cup\{e\}$ is a flat. Each summand contains
the factor $\tau(M/(F\cup\{e\}))$. In odd rank $s$, $\tau$ is the
coefficient of $t^{(s-1)/2}$ in the Kazhdan--Lusztig polynomial. In even
rank it is zero.

As in the proof of Theorem~3.2 of \cite{bv}, contracting $F\cup\{e\}$
leaves the two contracted sides meeting only in the vertex containing
$u$ and $v$. Their cycle matroids form a direct sum. Both summands have
positive rank. Indeed, at least one edge of each $H_i-e$ is missing from
$F$, since otherwise $F$ connects $u$ and $v$ and its closure contains $e$.
Each missing edge survives as a nonloop because $F\cup\{e\}$ is a flat.
Lemma~2.7 of \cite{bv} gives $\tau=0$ for a direct sum of two positive-rank
matroids. All correction terms therefore vanish.
\end{proof}

\begin{theorem}\label{polygon}
Let $G$ be the boundary of a convex polygon on $r+1\ge3$ vertices with
any set of pairwise noncrossing diagonals added. Then
\[
 Z_{M(G)}(t)=\sum_{j=0}^r \frac{1}{r+1}
        \binom{r+1}{j}\binom{r+1}{j+1}t^j=:N_{r+1}(t).
\]
\end{theorem}
\begin{proof}
Each diagonal separates the graph into two subpolygons satisfying
Lemma~\ref{lem}. Each remains connected after deleting the diagonal,
since it contains the corresponding boundary arc. Delete all diagonals
in turn. The remaining graph is the cycle on $r+1$ vertices, whose cycle
matroid is $U_{r,r+1}$. Its $Z$-polynomial is the displayed polynomial by
Proposition~4.9 of Proudfoot, Xu and Young \cite{pxy}.
\end{proof}

\section{Ribbons}
A ribbon is a connected skew diagram containing no $2\times2$ square.
Alexandersson asks whether the $Z$-polynomial of a connected ribbon
lattice path matroid depends only on the rank and equals $N_{r+1}$
\cite{alex}. The full fan case is the formula of Lu, Xie and Yang \cite{lxy}.

\begin{theorem}
Every connected ribbon lattice path matroid $M$ of rank $r$ has
$Z_M=N_{r+1}$. Its zeros are simple and strictly negative.
\end{theorem}
\begin{proof}
In rank one the lattice of flats is a chain on two elements, giving
$Z_M=1+t=N_2$. If rank-zero degenerate ribbons are admitted, their
lattice of flats has one element and $Z_M=1=N_1$.

Suppose $r\ge2$. The two bounding paths have no intermediate common
vertex. Otherwise the bases split as a direct sum on the nonempty sets
of step positions before and after that vertex. An interior lattice
point of the diagram would be incident to four unit squares of the
diagram, making a $2\times2$ square. Thus $M$ is a snake in the sense of
Knauer, Mart\'inez-Sandoval and Ram\'irez Alfons\'in \cite{kmr}. Their
Theorem~3.2 represents $M$ by a multi-fan, with a hub and the $r$
vertices of a path. Spokes join the hub to selected path vertices,
including both endpoints. The interval presentation gives one path
edge for each element shared by consecutive intervals, and a spoke
for each element belonging to a single interval.

Simplifying each parallel class preserves the ranked lattice of flats.
The path and its endpoint spokes form a polygonal boundary on $r+1$
vertices. The other spokes are noncrossing diagonals, so
Theorem~\ref{polygon} gives the identity. For every $r$, the Jacobi
binomial expansion gives
\[
 N_{r+1}(t)=\frac{(1-t)^r}{r+1}
 P_r^{(1,1)}\!\left(\frac{1+t}{1-t}\right).
\]
The $r$ simple Jacobi zeros in $(-1,1)$ map to negative real numbers
under $x\mapsto(x-1)/(x+1)$.
\end{proof}

\paragraph{Relation to earlier work.}
Braden and Vysogorets explicitly treat polygon dissections as iterated
parallel connections in their study of Kazhdan--Lusztig polynomials. The
same vanishing argument applied to their $Z$-deletion identity gives
Lemma~\ref{lem}. Ferroni, Matherne and Vecchi also record the vanishing
in their equivariant treatment \cite[Theorem~1.1 and \S4.1.2]{fmv}.
Our deduction identifies the consequence for ribbons. Neither its
validity nor its scope depends on a computation.

\small
\begin{thebibliography}{9}
\setlength{\itemsep}{2pt}
\bibitem{alex} P. Alexandersson, \emph{Real-rooted polynomials: open problems},
\href{https://www.symmetricfunctions.com/realRootedOpenProblems.htm\#openProblemsGraphs}{Symmetric Functions Catalog}, 2026, and \href{https://github.com/PerAlexandersson/RealRooted/issues/2650}{RealRooted issue 2650}, item~4, 24 November 2026.
\bibitem{bv} T. Braden and A. Vysogorets, Kazhdan--Lusztig polynomials of matroids under deletion, \emph{Electron. J. Combin.} 27 (2020), P1.17. \href{https://doi.org/10.37236/9026}{doi:10.37236/9026}.
\bibitem{fmv} L. Ferroni, J. P. Matherne, and L. Vecchi, Deletion formulas for equivariant Kazhdan--Lusztig polynomials of matroids, \emph{SIAM J. Discrete Math.} 39 (2025). \href{https://doi.org/10.1137/24M1673796}{doi:10.1137/24M1673796}.
\bibitem{kmr} K. Knauer, L. Mart\'inez-Sandoval, and J. L. Ram\'irez Alfons\'in, A Tutte polynomial inequality for lattice path matroids, \emph{Adv. Appl. Math.} 94 (2018), 23--38. \href{https://doi.org/10.1016/j.aam.2016.11.008}{doi:10.1016/j.aam.2016.11.008}.
\bibitem{lxy} L. Lu, M. H. Y. Xie, and A. L. B. Yang, Kazhdan--Lusztig polynomials of fan matroids, wheel matroids, and whirl matroids, \emph{J. Combin. Theory Ser. A} (2022). \href{https://doi.org/10.1016/j.jcta.2022.105665}{doi:10.1016/j.jcta.2022.105665}.
\bibitem{pxy} N. Proudfoot, Y. Xu, and B. Young, The $Z$-polynomial of a matroid, \emph{Electron. J. Combin.} 25 (2018), P1.26. \href{https://doi.org/10.37236/7105}{doi:10.37236/7105}.
\end{thebibliography}
\end{document}
